% GENERATED FILE. Edit canonical lessons, then run npm run generate:books.
% canonical-source-hash: fnv1a64:11c98cb550a34bac
% canonical-lessons: KA-C213-kote, KA-C213-kalyana-mantapa, KA-C213-jim, KA-C213-kridangana, KA-C213-vastusangrahalaya

\chapter{Fort, Wedding hall, Gym, Stadium, Museum}
\label{ch:ka-fort-wedding-hall-gym-stadium-museum}
\hlchaptermodality{213}

\begin{chapteropening}
\textbf{By the end of this chapter:} \emph{I can say a fort, a wedding hall, a gym, a stadium and a museum.}

Complete the last lesson of chapter 213: I can say a fort, a wedding hall, a gym, a stadium and a museum.
\end{chapteropening}

\section[\texorpdfstring{\kn{ಕೋಟೆ} (kōṭe)}{ಕೋಟೆ (kōṭe)}]{\kn{ಕೋಟೆ} (kōṭe) — a fort}

\label{lesson:KA-C213-kote}



\begin{quote}
Before the new one: say the Kannada for a ground, then the Kannada for a direction.
\end{quote}

\subsection*{You'll want to know: \kn{ಕೋಟೆ}}
\textbf{\kn{ಕೋಟೆ}} — \emph{kōṭe} — \textquotedblleft{}a fort\textquotedblright{}.

\begin{cousinweb}
\textbf{\kn{ಕೋಟೆ}} (\emph{kōṭe}) is native Dravidian, one word with Tamil \emph{kōṭṭai} (DEDR 2207). DEDR lists \emph{kōṭe} itself, \textquotedblleft{}fort, rampart\textquotedblright{}. DEDR 2207 compares Sanskrit \emph{koṭṭa}.
\end{cousinweb}

\begin{grammarlens}[title={What you've built}]
One more word for this chapter.
\end{grammarlens}

\subsection*{Your turn}
\begin{itemize}
\raggedright
  \item \emph{Say it:} \emph{kōṭe}
  \item \emph{Say it:} \emph{kōṭe}, once more
  \item \emph{Say it:} say \emph{kōṭe}
  \item \emph{Recall:} say the Kannada for a ground, then the Kannada for a direction, then say \emph{kōṭe} again
\end{itemize}

\begin{tcolorbox}[breakable,colback=teal!4,colframe=teal!35!black,title={Before you move on}]
What does \kn{ಕೋಟೆ} mean? (\textquotedblleft{}A fort\textquotedblright{}.) Say it once more.
\end{tcolorbox}
\section[\texorpdfstring{\kn{ಕಲ್ಯಾಣ} \kn{ಮಂಟಪ} (kalyāṇa maṇṭapa)}{ಕಲ್ಯಾಣ ಮಂಟಪ (kalyāṇa maṇṭapa)}]{\kn{ಕಲ್ಯಾಣ} \kn{ಮಂಟಪ} (kalyāṇa maṇṭapa) — a wedding hall}

\label{lesson:KA-C213-kalyana-mantapa}



\begin{quote}
Before the new one: say the Kannada for a direction, then the Kannada for a fort.
\end{quote}

\subsection*{You'll want to know: \kn{ಕಲ್ಯಾಣ} \kn{ಮಂಟಪ}}
\textbf{\kn{ಕಲ್ಯಾಣ} \kn{ಮಂಟಪ}} — \emph{kalyāṇa maṇṭapa} — \textquotedblleft{}a wedding hall\textquotedblright{}.

\begin{cousinweb}
\textbf{\kn{ಕಲ್ಯಾಣ} \kn{ಮಂಟಪ}} (\emph{kalyāṇa maṇṭapa}) is Sanskrit \emph{kalyāṇa}, \textquotedblleft{}beautiful, lucky\textquotedblright{}, the word for a wedding (CDIAL 2952), and \emph{maṇṭapa}, \textquotedblleft{}open pavilion\textquotedblright{} (CDIAL 9740): the pavilion for the wedding.
\end{cousinweb}

\begin{grammarlens}[title={What you've built}]
One more word for this chapter.
\end{grammarlens}

\subsection*{Your turn}
\begin{itemize}
\raggedright
  \item \emph{Say it:} \emph{kalyāṇa maṇṭapa}
  \item \emph{Say it:} \emph{kalyāṇa maṇṭapa}, once more
  \item \emph{Say it:} say \emph{kalyāṇa maṇṭapa}
  \item \emph{Recall:} say the Kannada for a direction, then the Kannada for a fort, then say \emph{kalyāṇa maṇṭapa} again
\end{itemize}

\begin{tcolorbox}[breakable,colback=teal!4,colframe=teal!35!black,title={Before you move on}]
What does \kn{ಕಲ್ಯಾಣ} \kn{ಮಂಟಪ} mean? (\textquotedblleft{}A wedding hall\textquotedblright{}.) Say it once more.
\end{tcolorbox}
\section[\texorpdfstring{\kn{ಜಿಮ್} (jim)}{ಜಿಮ್ (jim)}]{\kn{ಜಿಮ್} (jim) — a gym}

\label{lesson:KA-C213-jim}



\begin{quote}
Before the new one: say the Kannada for a fort, then the Kannada for a wedding hall.
\end{quote}

\subsection*{You'll want to know: \kn{ಜಿಮ್}}
\textbf{\kn{ಜಿಮ್}} — \emph{jim} — \textquotedblleft{}a gym\textquotedblright{}.

\begin{cousinweb}
\textbf{\kn{ಜಿಮ್}} (\emph{jim}) is a loan from English \emph{gym}.
\end{cousinweb}

\begin{grammarlens}[title={What you've built}]
One more word for this chapter.
\end{grammarlens}

\subsection*{Your turn}
\begin{itemize}
\raggedright
  \item \emph{Say it:} \emph{jim}
  \item \emph{Say it:} \emph{jim}, once more
  \item \emph{Say it:} say \emph{jim}
  \item \emph{Recall:} say the Kannada for a fort, then the Kannada for a wedding hall, then say \emph{jim} again
\end{itemize}

\begin{tcolorbox}[breakable,colback=teal!4,colframe=teal!35!black,title={Before you move on}]
What does \kn{ಜಿಮ್} mean? (\textquotedblleft{}A gym\textquotedblright{}.) Say it once more.
\end{tcolorbox}
\section[\texorpdfstring{\kn{ಕ್ರೀಡಾಂಗಣ} (krīḍāṅgaṇa)}{ಕ್ರೀಡಾಂಗಣ (krīḍāṅgaṇa)}]{\kn{ಕ್ರೀಡಾಂಗಣ} (krīḍāṅgaṇa) — a stadium}

\label{lesson:KA-C213-kridangana}



\begin{quote}
Before the new one: say the Kannada for a wedding hall, then the Kannada for a gym.
\end{quote}

\subsection*{You'll want to know: \kn{ಕ್ರೀಡಾಂಗಣ}}
\textbf{\kn{ಕ್ರೀಡಾಂಗಣ}} — \emph{krīḍāṅgaṇa} — \textquotedblleft{}a stadium\textquotedblright{}.

\begin{cousinweb}
\textbf{\kn{ಕ್ರೀಡಾಂಗಣ}} (\emph{krīḍāṅgaṇa}) is Sanskrit \emph{krīḍā}, \textquotedblleft{}play, sport\textquotedblright{}, and \emph{aṅgaṇa}, \textquotedblleft{}courtyard\textquotedblright{}: \textquotedblleft{}a yard for sport\textquotedblright{}.
\end{cousinweb}

\begin{grammarlens}[title={What you've built}]
One more word for this chapter.
\end{grammarlens}

\subsection*{Your turn}
\begin{itemize}
\raggedright
  \item \emph{Say it:} \emph{krīḍāṅgaṇa}
  \item \emph{Say it:} \emph{krīḍāṅgaṇa}, once more
  \item \emph{Say it:} say \emph{krīḍāṅgaṇa}
  \item \emph{Recall:} say the Kannada for a wedding hall, then the Kannada for a gym, then say \emph{krīḍāṅgaṇa} again
\end{itemize}

\begin{tcolorbox}[breakable,colback=teal!4,colframe=teal!35!black,title={Before you move on}]
What does \kn{ಕ್ರೀಡಾಂಗಣ} mean? (\textquotedblleft{}A stadium\textquotedblright{}.) Say it once more.
\end{tcolorbox}
\section[\texorpdfstring{\kn{ವಸ್ತುಸಂಗ್ರಹಾಲಯ} (vastusaṅgrahālaya)}{ವಸ್ತುಸಂಗ್ರಹಾಲಯ (vastusaṅgrahālaya)}]{\kn{ವಸ್ತುಸಂಗ್ರಹಾಲಯ} (vastusaṅgrahālaya) — a museum}

\label{lesson:KA-C213-vastusangrahalaya}



\begin{quote}
Before the new one: say the Kannada for a gym, then the Kannada for a stadium.
\end{quote}

\subsection*{You'll want to know: \kn{ವಸ್ತುಸಂಗ್ರಹಾಲಯ}}
\textbf{\kn{ವಸ್ತುಸಂಗ್ರಹಾಲಯ}} — \emph{vastusaṅgrahālaya} — \textquotedblleft{}a museum\textquotedblright{}.

\begin{cousinweb}
\textbf{\kn{ವಸ್ತುಸಂಗ್ರಹಾಲಯ}} (\emph{vastusaṅgrahālaya}) is three Sanskrit words: \emph{vastu}, \textquotedblleft{}thing\textquotedblright{} (CDIAL 11449), \emph{saṅgraha}, \textquotedblleft{}collection\textquotedblright{} (CDIAL 12852), and \emph{ālaya}, \textquotedblleft{}house\textquotedblright{} (CDIAL 1366): \textquotedblleft{}a house for a collection of things\textquotedblright{}.
\end{cousinweb}

\begin{grammarlens}[title={What you've built}]
One more word for this chapter.
\end{grammarlens}

\subsection*{Your turn}
\begin{itemize}
\raggedright
  \item \emph{Say it:} \emph{vastusaṅgrahālaya}
  \item \emph{Say it:} \emph{vastusaṅgrahālaya}, once more
  \item \emph{Say it:} say \emph{vastusaṅgrahālaya}
  \item \emph{Recall:} say the Kannada for a gym, then the Kannada for a stadium, then say \emph{vastusaṅgrahālaya} again
\end{itemize}

\begin{tcolorbox}[breakable,colback=teal!4,colframe=teal!35!black,title={Before you move on}]
What does \kn{ವಸ್ತುಸಂಗ್ರಹಾಲಯ} mean? (\textquotedblleft{}A museum\textquotedblright{}.) Say it once more.
\end{tcolorbox}
